\section{Introduction}
\label{sec:introduction}

The concept that physical reality is fundamentally constrained by information processing costs is rooted in Landauer's principle \cite{landauer_irreversibility_1961} and Wheeler's ``It from Bit'' hypothesis \cite{wheeler_information_1989}. More recently, Verlinde \cite{verlinde_origin_2011} and Bianconi \cite{bianconi_thermodynamics_2025} have proposed that gravity and spacetime are entropic phenomena, emerging from information gradients and quantum-entropy coupling. These approaches motivate a systems-first question: if physical structure is informationally constrained, what minimal architecture makes that information processing possible?

Existing programs such as Causal Set Theory \cite{bombelli_space-time_1987}, Loop Quantum Gravity
\cite{rovelli_quantum_2004}, and emergent gravity \cite{verlinde_origin_2011,jacobson_thermodynamics_1995} treat spacetime or gravity as fundamental starting points. This work instead tests deriving the physical structure using persistence mechanics.

\textit{Persistence Mechanics}\cite{kate_meyer_persistence_2026} (PM) formalizes the exact structural requirements for any system to exist over time. It establishes that maintaining coherence against entropic erasure requires solving three existential problems: Finiteness (bounded resources), Conservation (preserved information), and Causality (well-defined temporal evolution). These exact a strictly positive entropic cost governed by six irreducible pillars: Capacity, Map, Protocol, Governor, Toll, and Margin. 

While these constraints intuitively govern the survival of macroscopic complex systems (such as computational networks or biological cells), we pose a more fundamental question: how far down does this structural architecture apply? Does it stop at atoms, or does it dictate the laws of quantum physics? 

Here, we put this to the test. We apply the mechanics of persistence to the most demanding domain possible: the quantum vacuum. By treating the vacuum as a maximally constrained, zero-flux closed system—possessing no external memory, no deeper substrate beneath it, and no external boundary conditions. We ask if the sheer requirement of \textit{existence} uniquely determines its geometric structure. If persistence requirements alone cannot derive the vacuum's structure, then physical reality is governed by external arbitrary parameters. If they can, then all of physics can be understood as the limiting baseline of persistence mechanics.

We refer to this specific instantiation of PM applied to the vacuum as the \textbf{$E_8$-Persistence Theory}. We do not replace standard Effective Field Theory (EFT), but rather establish an information-theoretic dual. EFT provides a top-down continuum description of field dynamics; $E_8$-Persistence identifies the bottom-up discrete capacity constraints of the underlying geometric substrate from which those fields must emerge.


\subsection{The Minimal Vacuum Derivation}
We apply the six pillars of Persistence Mechanics to the vacuum in three strict, sequential layers, followed by empirical validation. 

\begin{enumerate}
    \item \textbf{System 0: Requirements.} We derive the geometric implications of Finiteness, Conservation, and Causality for the vacuum substrate under a strict Closure Constraint (zero external information) (\Cref{sec:System0_vacuum_requirements}).
    
    \item \textbf{System 1: Substrate.} System 0's requirements uniquely specify an $E_8$ lattice projected to $D=4$ spacetime. This yields unique structural invariants: a temporal cycle ($\Delta$), a discrete channel capacity ($\nu$), an interaction embedding dimension ($\sigma$), and a boundary topological characteristic ($\chi$) required by the architecture (\Cref{sec:system1_projection}).

    \item \textbf{System 2: Impedance.} We derive the geometric impedance ($Z_{\text{geo}}$) from the six-pillar entropic balance, quantifying the cost for a topological defect to persist on this substrate (\Cref{sec:system2_geometric_impedance}).

    \item \textbf{Empirical Correspondences.} We compare $Z_{\text{geo}}$ against experimental measurements of the fine-structure constant ($\alpha^{-1}$), identify a predicted Quantization Gap between kinematic and high-loop QED extractions, propose a falsification criterion, and derive the geometric origins of elementary charge ($e$) and the von Klitzing constant ($R_K$) (\Cref{sec:empirical_correspondences}).
\end{enumerate}

This paper is restricted to the static, quiescent-equilibrium baseline. Dynamics, Lagrangians, gauge structure, and ultraviolet Lorentz invariance restoration represent natural extensions, but they are not prerequisites for the architectural test performed here: establishing whether the universal constraints of persistence uniquely determine the vacuum substrate.

To contextualize this framework against prior unification attempts, the Distler--Garibaldi no-go theorem~\cite{distler_there_2010} applies strictly to continuous $E_8$ gauge algebras. Because the $E_8$ substrate established in System 1 is a discrete geometric lattice, it intrinsically resolves this constraint by isolating mirror degrees of freedom outside the causal manifold (\Cref{app:geometric_origin_of_chiral_fermions}).

A critical attribute of this methodology must be explicitly emphasized: The derived structural invariants are strict, mathematical consequences of applying eliminative persistence constraints to a discrete substrate. Because the substrate admits zero continuous degrees of freedom, the resulting geometric impedance $Z_{geo}$ is a strictly parameter-free prediction.